\begin{abstract}
Large-scale unsupervised language models (LMs) acquire extensive world knowledge and display some reasoning abilities, but their inherently unsupervised training regime makes it challenging to precisely control their outputs. Current approaches to enhance model controllability rely on gathering human judgments about the relative merit of generated outputs, and subsequently adjusting the model through fine-tuning so that its responses are more in line with human preferences; this is frequently done via reinforcement learning from human feedback (RLHF). However, the RLHF framework involves multiple complicated and potentially unstable steps: first, constructing a reward model that encodes the preferences gathered from humans, and then applying reinforcement learning to further train the base LM to optimize this learned reward, while ensuring that the model does not deviate excessively from its pretrained state. In this work, we propose an alternative parameterization for the reward model in RLHF, which makes it possible to derive the associated optimal policy in closed form. This insight reduces the original RLHF objective to a straightforward classification loss, leading to a novel approach called \textit{Direct Preference Optimization} (DPO). DPO offers a stable and efficient solution that is easy to compute, as it avoids the need for sampling from the LM during fine-tuning and reduces the burden of extensive hyperparameter selection. Our empirical results indicate that DPO fine-tunes LMs to follow human preferences as effectively as—or better than—prior techniques. In particular, DPO not only surpasses PPO-based RLHF in controlling generation sentiment, but also achieves comparable or superior performance in summarization and single-turn dialogue tasks, all while remaining much easier to implement and train.
\end{abstract}

\section{Introduction}
Large-scale unsupervised language models (LMs), when trained on expansive datasets, have demonstrated unexpectedly advanced abilities~\citep{chowdhery2022palm, brown2020language, touvron2023llama,bubeck2023sparks}. Despite these strengths, their training data comprises human-generated content marked by diverse objectives, levels of expertise, and varying intent. Many of these underlying motivations and skills are not always desirable as targets for imitation; for instance, although an AI programming assistant should be capable of \textit{recognizing} common coding errors to address them, we ultimately want the model to prioritize the higher-quality coding practices—even if such examples are rare in its corpus—during code generation. Analogously, even if a language model must be \textit{informed} about a popular misconception accepted by half the population, it should avoid perpetuating that misconception in its responses at an equivalent rate. This distinction between a model’s extensive \textit{knowledge and skills} and the subset of \emph{desired responses and behaviors} we wish it to manifest is essential for building AI systems that are reliable, high-performing, and easily governed \citep{ouyang2022training}. Though current practice often involves aligning LMs with human values using reinforcement learning (RL) approaches, we will demonstrate that the RL-based objectives employed in prevalent methods can, in fact, be directly optimized via a straightforward binary cross-entropy objective. This offers a substantial simplification of the preference learning workflow.

Generally, present techniques impart desired characteristics to language models by utilizing carefully constructed sets of human preference data that embody the qualities deemed safe and useful by users. This process of preference learning is performed subsequent to an initial stage of extensive unsupervised pre-training on text corpora. Although supervised fine-tuning with human-curated, high-quality response examples is the most direct method of preference learning, the most effective approaches currently are based on reinforcement learning from human or AI feedback (RLHF/RLAIF; \citep{christiano2017deep,bai2022constitutional}). RLHF methodologies construct a reward model from human preference data and then employ RL to update the language model policy to favor outputs with high reward, while restraining deviation from the base model. Despite RLHF’s success in achieving strong conversational and programming performance, its implementation introduces considerable complexity compared to supervised fine-tuning, as it necessitates training several LMs and performing policy sampling within the optimization loop—substantially increasing computational overhead.

In this work, we demonstrate a method for directly tuning a language model to align with human preferences, bypassing the need for explicit reward modeling or reinforcement learning. We introduce \textit{{Direct Preference Optimization} (DPO)}, an algorithm that, while implicitly optimizing the same objective as established RLHF approaches (namely, reward maximization subject to a KL-divergence regularization), offers a simpler implementation and easier training dynamics. At a high level, DPO updates the model by amplifying the log probability of favored responses relative to those less preferred, utilizing an adaptive importance weighting mechanism at the individual example level to mitigate the model collapse observed with naive probability ratio-based objectives. Similar to prior techniques, DPO leverages a theoretical preference framework (such as the Bradley-Terry model; \cite{bradley1952rankanalysis}) to quantify alignment between a reward signal and observed preference data. In contrast to traditional pipelines—which employ the preference model to craft a loss for training a reward predictor, followed by policy optimization on that learned reward—DPO re-expresses the preference loss directly in terms of the policy parameters through a variable transformation. Provided with a set of human preference annotations comparing model responses, DPO enables policy optimization through a straightforward binary cross-entropy loss, leading to an optimal policy corresponding to a latent reward consistent with the empirical preference judgments.

The central contribution of this work is the Direct Preference Optimization (DPO) method, a straightforward and RL-free approach for preference-based training of language models. Empirical results indicate that DPO matches or surpasses existing techniques, including RLHF variants based on PPO, in preference learning benchmarks such as sentiment control, summarization, and dialogue, using models with up to 6B parameters.

As self-supervised language models have scaled up, they have demonstrated the capacity to solve certain tasks in a zero-shot fashion \citep{radford2019language} or with limited examples provided in the prompt (few-shot learning) \citep{gpt3,megatron,chowdhery2022palm}. Nevertheless, their efficacy on specific downstream applications and the degree to which their outputs align with user intentions are markedly enhanced when they undergo fine-tuning using datasets composed of instructions and human-generated completions \citep{mishra-etal-2022-cross,sanh2022multitask,chung2022scaling,thoppilan2022lamda}. This process, referred to as `instruction-tuning', enables large language models (LLMs) to extrapolate to novel instructions beyond those present in the instruction-tuning corpus and generally renders them more accessible and user-friendly \citep{chung2022scaling}. While instruction tuning has achieved considerable success, gathering \textit{relative} human assessments of answer quality typically proves less demanding than obtaining expert demonstrations. Consequently, subsequent research has adopted the approach of further fine-tuning LLMs with datasets constructed from human preference judgments, resulting in improvements across tasks such as translation \citep{kreutzer-etal-2018-reliability}, summarization \citep{stiennon2022learning,ziegler2020finetuning}, creative writing \citep{ziegler2020finetuning}, and instruction adherence \citep{ouyang2022training,ramamurthy2023is}. In these frameworks, a reward function, represented by a neural network, is initially trained to match the dataset of human preferences, typically utilizing models like the Bradley-Terry framework \citep{bradley1952rankanalysis}. Subsequently, the language model is further optimized to maximize these rewards through reinforcement learning techniques such as REINFORCE \citep{williams1992reinforce}, proximal policy optimization (PPO; \cite{schulman2017proximal}), or their extensions \citep{ramamurthy2023is}. Related work has involved employing LLMs that have already been instruction-tuned with human feedback to autonomously generate synthetic preference datasets focused on criteria like safety and harmlessness \citep{bai2022constitutional}. In these approaches, human supervision is limited to supplying textual guidelines that inform the LLM’s annotation process. Collectively, these advances fuse research from two key domains: one concerned with applying reinforcement learning to train language models for diverse objectives~\citep{Ranzato2015SequenceLT,paulus2018a,wu2018learning}, and another centered on frameworks for learning from human preference data \citep{christiano2017deep,kupcsik2018learning}. Nonetheless, despite the attractiveness of training with comparative human preferences, the practicalities of fine-tuning large-scale models using reinforcement learning introduce considerable complexity; the present work contributes a methodologically rigorous alternative for optimizing models based on relative preferences, circumventing the need for reinforcement learning.

Beyond linguistic applications, the problem of learning policies based on preferences has been explored extensively within both bandit frameworks and reinforcement learning, with a range of methodologies proposed. In particular, contextual bandit algorithms that utilize preferences or rankings in place of explicit reward signals are referred to as contextual dueling bandits (CDB; \cite{yue2012karmed,dudik2015contextual}). In situations where absolute reward feedback is unavailable, analyses of CDBs employ the concept of a \textit{von Neumann winner}—a policy that, on average, wins against any competing policy at least half the time \citep{dudik2015contextual}. A notable distinction, however, is that CDBs assume access to online preference feedback, while approaches focused on human preferences typically work with a static collection of offline, annotated preference pairs \citep{yan2022human}. Analogously, \textit{preference-based RL} (PbRL) methods utilize binary preferences obtained via an \textit{unknown} scoring function, as opposed to direct reward signals \citep{BusaFekete2014,ruiz2023dueling}. Existing PbRL algorithms often require estimating a latent reward (or scoring) function from preference data—an explicit intermediate step—before optimizing with respect to that function \citep{jain2013learning,BusaFekete2014,christiano2017deep,sadigh2017active,kupcsik2018learning}, and can accommodate off-policy preference datasets. In contrast, we propose a unified framework for policy learning that dispenses with the intermediate reward modeling step, directly optimizing the policy with respect to observed preferences.

\section{Preliminaries}\label{section:prelims}

We outline the RLHF pipeline introduced by \citeauthor{ziegler2020finetuning} (with extensions in \citep{stiennon2022learning, bai2022training, ouyang2022training}), which is conventionally organized into three distinct stages: (1) supervised fine-tuning (SFT); (2) collection of preferences and reward model training; and (3) reinforcement learning.

\textbf{SFT}: The initial stage in RLHF involves adapting a pre-trained language model to the specific target task (such as dialogue or summarization) through supervised training on high-quality, task-relevant data, resulting in a model $\pi^\text{SFT}$.

\textbf{Reward Modelling Phase}: In this stage, the SFT-adapted model is given prompts $x$ and produces candidate responses $(y_1, y_2)\sim \pi^\text{SFT}(y \mid x)$. Human annotators are then asked to express their preference for one response over the other, formalized as $y_w\succ y_l \mid x$—here, $y_w$ represents the selected (preferred) output and $y_l$ the rejected one, for each prompt-response pair. The preference mechanism is hypothesized to result from an underlying, unobserved reward function $r^*(y, x)$. Multiple statistical models exist for representing such preference data, with the Bradley-Terry (BT) model \cite{bradley1952rankanalysis} being a common option (alternatively, more flexible models such as Plackett-Luce \citep{plackett1975analysis, luce2012individual} are also suitable if multiple rankings per prompt are available). According to the BT framework, the empirical human preference distribution $p^*$ is defined as:

Given a static dataset of pairwise comparisons $\mathcal{D}=\bigl\{x^{(i)}, y_w^{(i)}, y_l^{(i)}\bigr\}_{i=1}^N$ sampled from the distribution $p^*$, one can define a reward function $r_{\phi}(x, y)$ and learn its parameters via the principle of maximum likelihood. Treating the problem as binary classification, the negative log-likelihood objective can be written as:

\begin{equation}\label{eq:reward_model}
    \mathcal{L}_R(r_{\phi}, \mathcal{D}) = -\mathbb{E}_{(x, y_w, y_l)\sim \mathcal{D}}\bigl[\log \sigma(r_{\phi}(x, y_w)- r_{\phi}(x, y_l))\bigr]
\end{equation}

where $\sigma$ denotes the logistic sigmoid function. In language modeling, $r_{\phi}(x, y)$ is typically initialized from the SFT policy $\pi^\text{SFT}(y \mid x)$, with an additional linear head placed atop the final transformer layer to output a scalar reward estimate \cite{ziegler2020finetuning}. To control reward variability, existing studies standardize the outputs so that $\mathbb{E}_{x,y\sim \mathcal{D}}\left[r_\phi(x, y)\right] = 0$ for any input $x$.

\textbf{RL Fine-Tuning Phase}: Once the reward model is learned, it serves as a feedback mechanism in the subsequent RL phase. Following conventions established in previous work~\citep{jaques2017sequence, jaques2020human}, the optimization problem is expressed as

\begin{equation}\label{eq:RL}
\max_{\pi_{\theta}}  \mathbb{E}_{x\sim \mathcal{D}, y\sim \pi_{\theta}(y \mid x)}\bigl[r_{\phi}(x, y)\bigr] - \beta\mathbb{D}_{\textrm{KL}}\bigl[\pi_{\theta}(y\mid x)\mid \mid \pi_\text{ref}(y\mid x)\bigr],
\end{equation}

with $\beta$ controlling the regularization towards the reference policy $\pi_\text{ref}$, typically the initial SFT model $\pi^\text{SFT}$. The language model policy $\pi_\theta$ is commonly initialized from $\pi^\text{SFT}$. This regularization term is crucial as it constrains the policy within regions where the reward model provides reliable guidance, maintains sample diversity, and prevents degeneration to singular high-reward outputs. Owing to the discrete nature of text generation, direct optimization of this objective is infeasible via gradient-based methods, thus necessitating reinforcement learning approaches. Standard practice \citep{ziegler2020finetuning, stiennon2022learning, bai2022training, ouyang2022training} is to define the reward as ${r(x, y) = r_{\phi}(x, y) -\beta (\log \pi_{\theta}(y\mid x) - \log \pi_\text{ref}(y\mid x))}$, and employ PPO \cite{schulman2017proximal} for optimization.

\section{Direct Preference Optimization}\label{sec:DPO}

Recognizing the difficulties encountered when applying reinforcement learning techniques to large models, such as those used in language modeling, we propose a method that directly optimizes policies based on preference data. Distinct from earlier RLHF strategies, which decouple the reward learning and RL optimization steps, our approach introduces a particular form of reward function that admits an explicit, closed-form solution for the associated optimal policy—bypassing the need for a reinforcement learning training loop. The central insight involves utilizing an analytic correspondence between the reward model and the optimal policy; this facilitates reformulating a loss defined over rewards as one over policies. Consequently, our approach circumvents the need to fit a separate reward network, while still conforming to established preference-based frameworks such as the Bradley-Terry model. Essentially, the policy network itself is structured to encapsulate both the behavioral policy and its underlying (implicit) reward.

\textbf{Deducing the DPO objective.} Beginning with the standard reinforcement learning (RL) objective from earlier work, as described in Eq.~\ref{eq:RL}, under an arbitrary reward function $r$, we follow the methodology of~\citet{peters2007reinforcement, peng2019advantage, korbak2022reinforcement, go2023aligning}. One can directly verify that the solution optimizing the KL-regularized reward maximization objective in Eq.~\ref{eq:RL} admits the closed-form expression:

\begin{equation}\label{eq:op_policy}
    \pi_r(y\mid x) = \frac{1}{Z(x)}\pi_\text{ref}(y\mid x)\exp\left(\frac{1}{\beta}r(x, y)\right),
\end{equation}

with the normalization factor $Z(x) =\sum_{y}\pi_\text{ref}(y\mid x)\exp\left(\frac{1}{\beta}r(x, y)\right)$ termed the partition function. For a full derivation, consult Appendix \ref{app:derivation1}. Even after replacing $r$ with the MLE-based reward estimator $r_{\phi}$ for the underlying ground-truth $r^*$, estimating $Z(x)$ remains computationally burdensome~\citep{korbak2022reinforcement, go2023aligning}, thereby complicating practical use of this representation. Nevertheless, by manipulating Eq.~\ref{eq:op_policy}, we can reformulate the reward as a function of the induced optimal policy $\pi_r$, the baseline policy $\pi_\text{ref}$, and the partition function $Z(\cdot)$. This is done by applying the logarithm to both sides of Eq.~\ref{eq:op_policy} and rearranging, yielding:

\begin{equation}\label{eq:main_eq}
    r(x,y) =\beta \log \frac{\pi_r(y\mid x)}{\pi_\text{ref}(y\mid x)} + \beta \log Z(x).
\end{equation}

Applying this form to both the true reward $r^*$ and its optimal associated policy $\pi^*$, we leverage that the Bradley-Terry model considers only reward differences between outcome pairs: ${p^*(y_1 \succ y_2 \mid x) = \sigma(r^*(x, y_1) - r^*(x, y_2))}$. Plugging the expression for $r^*(x,y)$ from Eq.~\ref{eq:main_eq} into the preference model of Eq.~\ref{eq:bradley-terry} eliminates the partition function, such that the resulting human preference likelihood depends exclusively on $\pi^*$ and the reference policy $\pi_\text{ref}$. As a result, the optimal RLHF policy $\pi^*$, as modeled under the Bradley-Terry framework, obeys:

\begin{equation}\label{eq:objective}
    p^*(y_1\succ y_2 \mid x)=\frac{1}{1 + \exp\left(\beta \log \frac{\pi^*(y_2\mid x)}{\pi_\text{ref}(y_2\mid x)} - \beta \log \frac{\pi^*(y_1\mid x)}{\pi_\text{ref}(y_1\mid x)}\right)}
\end{equation}

Derivational details can be found in Appendix~\ref{app:derivation2}. While Eq.~\ref{eq:objective} specifically involves the Bradley-Terry likelihood, a parallel argument can be made for more general Plackett-Luce models~\citep{plackett1975analysis, luce2012individual}, as outlined in Appendix~\ref{app:plackett_luce_models}.

This formulation, expressing human preference data probabilities in terms of the optimal policy rather than the reward parameterization, permits defining a maximum likelihood training criterion for a parameterized policy $\pi_\theta$. Mirroring the reward model-based strategy (see Eq.~\ref{eq:reward_model}), we express our objective for policy learning as:

\begin{equation}\label{eq:optimum_model}
    \mathcal{L}_\text{DPO}(\pi_{\theta}; \pi_\text{ref}) = -\mathbb{E}_{(x, y_w, y_l)\sim \mathcal{D}}\left[\log \sigma \left(\beta \log \frac{\pi_{\theta}(y_w\mid x)}{\pi_\text{ref}(y_w\mid x)} - \beta \log

In this approach, we represent the implicit reward with an alternative parameterization, where the resulting optimal policy coincides with $\pi_\theta$. Additionally, because this procedure corresponds to optimizing a reparameterized Bradley-Terry model, it inherits desirable theoretical characteristics—such as consistency, provided appropriate conditions on the distribution of preference data are satisfied \cite{bong2022generalized}. We elaborate further on the theoretical aspects of DPO and its relationship to existing methods in Section~\ref{sec:theory}.

\textbf{What mechanism underlies the DPO update?} To shed light on how DPO operates, it is helpful to inspect the gradient of the loss function $\mathcal{L}_\text{DPO}$. The derivative with respect to the parameters $\theta$ can be expressed as:

\begin{multline*}\label{eq:gradient}
    \nabla_\theta \mathcal{L}_\text{DPO}(\pi_\theta;\pi_\text{ref}) = \\ -\beta\mathbb{E}_{(x, y_w, y_l) \sim \mathcal{D}} \bigg[\underbrace{\sigma(\hat{r}_\theta(x, y_l) - \hat{r}_\theta (x, y_w))}_\text{assigns more weight if the reward estimate is incorrect}\bigg[\underbrace{\nabla_\theta\log \pi(y_w \mid x)}_\text{boosts likelihood of $y_w$} - \underbrace{\nabla_\theta\log\pi(y_l \mid x)}_\text{suppresses likelihood of $y_l$}\bigg]\bigg],
\end{multline*}

where $\hat{r}_\theta(x, y) = \beta \log \frac{\pi_\theta(y \mid x)}{\pi_\text{ref}(y \mid x)}$ defines the reward assigned by the language model $\pi_\theta$ relative to the reference $\pi_\text{ref}$ (see Section~\ref{sec:theory} for details). The interpretation of the gradient for $\mathcal{L}_\text{DPO}$ is that it promotes higher likelihood for the completions $y_w$ that are preferred and lowers the likelihood for those $y_l$ that are not. Crucially, the influence of each example is modulated by the extent to which the implicit reward model $\hat{r}_\theta$ overestimates the disfavored responses, scaled by $\beta$; that is, the adjustment reflects the degree of error in the ranking of completions, with $\beta$ determining the effect of the KL regularization. Empirical findings highlight that this weighting is critical—if one omits the weighting term, a straightforward variant of this approach can cause the language model to deteriorate in quality (refer to Appendix Table~\ref{tab:unlikelihood_generations}).

\textbf{DPO pipeline summary.} The typical procedure for DPO consists of: 1) Drawing completion samples $y_1, y_2 \sim \pi_\text{ref}(\cdot \mid x)$ for each prompt $x$, then assigning human preference labels to create the offline preference dataset $\mathcal{D} = \{x^{(i)}, y_w^{(i)}, y_l^{(i)}\}_{i=1}^N$; and 2) training the language model $\pi_\theta$ by minimizing $\mathcal{L}_\text{DPO}$, with respect to the given $\pi_\text{ref}$, dataset $\mathcal{D}$, and a selected value of $\beta$. In practice, practitioners often prefer to make use of existing public preference datasets instead of collecting new human preference annotations. Since these datasets are typically generated with $\pi^\text{SFT}$, we set $\pi_\text{ref} = \pi^\text{SFT}$ when this model is available. Conversely, if $\pi^\text{SFT}$ is absent, we set $\pi_\text{ref}$ by fitting it to the preferred completions ${(x, y_w)}$, formally: ${\pi_\text{ref} = \argmax_{\pi}\mathbb{E}_{x, y_w \sim \mathcal{D}}\left[\log \pi(y_w \mid x)\right]}$. This strategy addresses the issue of distribution shift between the unobserved true reference distribution and the one adopted by DPO, $\pi_\text{ref}$. Additional implementation specifics and hyperparameter choices are described in Appendix~\ref{app:implementation}.

\section{Theoretical Analysis of DPO} This section presents additional theoretical perspectives on the DPO approach, provides a formal justification for its use, and explains how DPO resolves problems seen in actor-critic based RLHF algorithms (notably PPO~\cite{schulman2017proximal}).

\label{sec:theory}

\subsection{Your Language Model Is Secretly a Reward Model} DPO eliminates the need to separately fit a reward model or conduct RL to train the policy by relying entirely on a maximum likelihood training objective. Observe that the optimization in Eq. \ref{eq:main_eq} corresponds to a Bradley-Terry model in which rewards are given by $r^*(x, y) = \beta \log\frac{\pi^*_\theta(y \mid x)}{\pi_\text{ref}(y \mid x)}$. Training the parameterized model $\pi_{\theta}$ is therefore mathematically equivalent to optimizing the reward model, as described in Eq. \ref{eq:reward_model}, under an appropriate variable transformation. In the analysis below, we formalize this reparameterization, demonstrate that it does not restrict the set of representable reward models, and prove that it can lead to the recovery of the optimal policy. We begin by establishing an equivalence relation among reward functions.

\begin{definition}
We define two reward functions $r(x, y)$ and $r'(x, y)$ as equivalent if and only if ${r(x, y)-r'(x, y) = f(x)}$ for some function $f$.
\end{definition}

This definition establishes an equivalence relation, dividing the collection of reward functions into distinct equivalence classes. The following lemmas can be derived:

\begin{lemma}\label{lemma:same_prefrence}
Within the Plackett-Luce, and specifically the Bradley-Terry, preference model, any two reward functions belonging to the same equivalence class yield identical preference distributions.
\end{lemma}

\begin{lemma}\label{lemma:same_policy}
    Any two reward functions within the same equivalence class will result in the same optimal policy for the associated constrained RL problem.
\end{lemma}

Since the proofs are elementary, we refer the reader to Appendix \ref{app:lemma1} for details. The first lemma reflects a well-documented under-specification property of models in the Plackett-Luce family \cite{plackett1975analysis}. Because of this issue, identifiability constraints are commonly imposed in order to ensure meaningful MLE solutions to Eq. \ref{eq:reward_model} \cite{bong2022generalized}. The second lemma indicates that all reward functions in an equivalence class prescribe the same optimal policy. Consequently, for our purposes, it suffices to identify any representative reward function from the relevant equivalence class. The following theorem, whose proof appears in Appendix~\ref{app:thm1}, formalizes this point:

\begin{theorem}\label{thm:main}
    Given mild conditions, every reward equivalence class consistent with Plackett-Luce (and specifically Bradley-Terry) models can be characterized by the reparameterization ${r(x, y) = \beta \log \frac{\pi(y\mid x)}{\pi_\text{ref}(y\mid x)}}$ for some model $\pi(y\mid x)$ and a fixed reference model $\pi_\text{ref}(y \mid x)$.
\end{theorem}

\begin{sproof}
    Take any reward function $r(x, y)$, with its corresponding optimal policy $\pi_r(y \mid x)$ defined as in Eq. \ref{eq:op_policy}. We seek to express a function from the equivalence class of $r$ via the proposed reparameterization. Define the projection
\begin{equation}
    f(r; \pi_\text{ref}, \beta)(x, y) = r(x, y) - \beta\log\sum_{y}\pi_\text{ref}(y\mid x)\exp\left(\frac{1}{\beta}r(x, y)\right)
\end{equation}
Here, $f$ normalizes the reward by subtracting the logarithm of the partition function of $\pi_r$. The normalization term is solely a function of $x$, so $f(r; \pi_\text{ref}, \beta)(x, y)$ remains within the equivalence class defined by $r(x, y)$. Upon substituting $r$ using the right-hand side of Eq.~\ref{eq:main_eq}—which is applicable for all reward functions—we obtain $f(r; \pi_\text{ref}, \beta)(x, y) = \beta \log \frac{\pi_r(y\mid x)}{\pi_\text{ref}(y\mid x)}$. Thus, this projection yields an equivalent reward function of the requisite form, demonstrating that the proposed reparameterization captures the full generality of the model class.
\end{sproof}

An alternative interpretation of Theorem~\ref{thm:main} is that it delineates precisely which reward function within each equivalence class is chosen by the DPO reparameterization. Specifically, it selects the reward function that fulfills:

\begin{equation}\label{eq:lag_p}
     \sum_{y}\underbrace{\pi_\text{ref}(y\mid x)\exp\left(\frac{1}{\beta}r(x, y)\right)}_{=\pi(y\mid x)\text{, by Thm.~\ref{thm:main} reparam.}} = 1,
\end{equation}

That is, $\pi(y\mid x)$ constitutes a legitimate probability distribution, where all probabilities are positive and total to 1. Examining Eq.~\ref{eq:op_policy} in conjunction with Eq.~\ref{eq:lag_p} reveals that Eq.~\ref{eq:lag_p} represents the partition function corresponding to the optimal policy derived from the reward function $r(x, y)$. The fundamental idea underlying the DPO approach is that we can introduce specific constraints into the generally underdetermined Plackett-Luce (notably including the Bradley-Terry model) class of preference models. By doing so, we retain the expressive capacity of representable reward models, while ensuring that the optimal policy defined in Eq.~\ref{eq:op_policy} remains analytically manageable for every prompt $x$.

\subsection{Instability of Actor-Critic Algorithms} Our framework is also useful for highlighting issues related to the instability commonly encountered with conventional actor-critic algorithms such as PPO, when applied to RLHF. Adhering to the RLHF methodology, our attention is on the RL fine-tuning phase as described in Section \ref{section:prelims}. We can establish parallels to the control as inference perspective \cite{levine2018reinforcement} applied to the constrained RL formulation in \ref{eq:RL}. Assuming a parameterized policy $\pi_{\theta}(y\mid x)$, the goal becomes minimizing $\mathbb{D}_{\text{KL}}[\pi_{\theta}(y|x) \mid \mid \pi^*(y\mid x)]$, where $\pi^*$ denotes the optimal policy from Eq.~\ref{eq:optimum_model} determined by the reward $r_{\phi}(y, x)$. A straightforward algebraic manipulation yields the following optimization problem:

\begin{equation}\label{eq:AC}
    \max_{\pi_{\theta}}\mathbb{E}_{\pi_{\theta}(y\mid x)}\bigg[\underbrace{r_{\phi}(x, y) -\beta\log\sum_{y}\pi_\text{ref}(y\mid x)\exp\left(\frac{1}{\beta}r_{\phi}(x, y)\right)}_{f(r_{\phi}, \pi_\text{ref}, \beta)} - \underbrace{\beta\log\frac{\pi_{\theta}(y\mid x)}{\pi_\text{ref}(y\mid x)}}_{\text{KL}}\bigg]
\end{equation}

The objective discussed here matches that optimized in earlier studies 
\citep{ziegler2020finetuning, stiennon2022learning, bai2022training, ouyang2022training} through a reward in the DPO-equivalent class for $r_{\phi}$. In this context, the normalization factor within $f(r_{\phi}, \pi_\text{ref}, \beta)$ can be viewed as representing the soft value function associated with the reference policy $\pi_\text{ref}$. Although this normalization does not modify the optimal policy itself, its omission may result in large-variance policy gradients, which can impede stable training. One can incorporate this normalization by introducing a learned value function, but optimizing such a function poses its own challenges. As an alternative, prior research has normalized rewards relative to a human completion baseline, which effectively serves as a Monte Carlo estimate—using a single sample—of the normalization term. By contrast, the DPO reparameterization derives a reward function inherently free of the need for baselines.

\section{Experiments}
This section empirically investigates DPO's effectiveness in policy training based purely on preference data. We begin by examining a controlled text generation problem to explore how DPO balances the objectives of reward maximization and minimization of KL-divergence from the reference policy, comparing its efficiency to widely used preference learning techniques like PPO. Subsequently, we analyze the scalability and robustness of DPO when applied to more complex models and challenging RLHF benchmarks, such as tasks involving summarization and dialogue. Our findings indicate that DPO consistently matches or outperforms strong baselines, such as PPO-based RLHF and best-of-$N$ sampling under a learned reward, while requiring minimal hyperparameter adjustment. Preceding these empirical results, we outline our experimental protocol, with supplementary information provided in Appendix~\ref{app:exp_details}.

\textbf{Tasks.} In our empirical studies, we consider three distinct open-ended text generation scenarios. Across all tasks, the methods under investigation are trained to optimize a policy using a dataset of preference comparisons $\mathcal{D}=\bigl\{x^{(i)}, y_w^{(i)}, y_l^{(i)}\bigr\}_{i=1}^N$. In the case of \textbf{controlled sentiment generation}, $x$ represents an initial segment from a movie review selected from the IMDb corpus \cite{maas-EtAl:2011:ACL-HLT2011}, and the model is required to continue the text such that $y$ conveys a positive sentiment. To rigorously assess model behavior, we construct preference pairs for generations with the aid of an external sentiment classifier, enforcing that $p(\text{positive}\mid x,y_w)>p(\text{positive}\mid x,y_l)$. For SFT, we conduct further training of GPT-2-large to convergence using the training portion of IMDb reviews (see App~\ref{app:sentiment_details} for more information). In the \textbf{summarization} setting, $x$ refers to a Reddit forum post, and the task is to generate a concise summary $y$ capturing its primary messages. Our approach is consistent with previous research, relying on the Reddit TL;DR summarization dataset \citep{volske-etal-2017-tl} in conjunction with a set of human preference annotations provided by \citeauthor{stiennon2022learning}. Here, we adopt an SFT model fine-tuned on summaries composed by humans,\footnote{\url{https://huggingface.co/CarperAI/openai_summarize_tldr_sft}} and perform RLHF with the TRLX \citep{leandro_von_werra_2023_7790115} library. The utilized preference data comes from \citeauthor{stiennon2022learning}, who collected annotations on outputs sampled from a related SFT system. Lastly, the \textbf{single-turn dialogue} task features $x$ as a user prompt—ranging from scientific queries to personal advice requests. The system is tasked with generating a reply $y$ that is both informative and engaging. For this, we leverage the Anthropic Helpful and Harmless dialogue collection \citep{bai2022training}, which consists of 170,000 exchanges between humans and an AI assistant. Each dialogue concludes with two alternative AI responses and an indicator reflecting human preference. As there is no established SFT model for this task, we construct one by fine-tuning a standard pretrained language model solely on responses preferred by annotators.

\textbf{Evaluation.} Our experimental evaluation relies on two complementary strategies. To measure how well each algorithm performs in maximizing rewards under constraints, we examine the relationship between reward achieved and KL-divergence from the reference policy in the controlled sentiment generation scenario. This analysis is enabled by our access to a ground-truth reward function—specifically, a sentiment classifier—which allows precise calculation of the reward frontier. However, such ground-truth rewards are typically unavailable outside this controlled context. As a result, for real-world tasks we assess algorithms by their \textit{win rate} against a baseline policy, employing GPT-4 as a surrogate for human judgment when rating the quality of summaries or the helpfulness of dialogue responses in the summarization and single-turn dialogue tasks, respectively. In summarization, the baseline consists of reference summaries from the test split; for dialogue, it is the response preferred by annotators in the test set. While prior research has shown that large language models can outperform standard metrics as automatic evaluators \citep{Chen2023ExploringTU}, we further substantiate our choice of GPT-4 through a dedicated human study described in Sec.~\ref{sec:human-judgments}. Our findings demonstrate that GPT-4's assessments exhibit high concordance with human judgments, with the level of agreement between humans and GPT-4 generally matching or exceeding agreement between human annotators themselves.

\textbf{Methods.} Alongside DPO, we benchmark multiple established approaches for aligning language models to human preferences. For the summarization task, we include zero-shot prompting with \textbf{GPT-J} \citep{gpt-j}, and for the dialogue task, we use 2-shot prompting with \textbf{Pythia-2.8B} \citep{biderman2023pythia}. We also assess the \textbf{SFT} model as well as \textbf{Preferred-FT}, which is generated by further fine-tuning on the chosen completion $y_w$, sourced from either the SFT model (for controlled sentiment and summarization tasks) or a general language model (for single-turn dialogue). Another relevant approach is \textbf{Unlikelihood}~\citep{welleck2019neural}, which trains the model to maximize the likelihood of $y_w$ while simultaneously \emph{minimizing} the likelihood of $y_l$; here, the influence of the unlikelihood loss can be scaled by a tunable coefficient $\alpha\in[0,1]$. We further compare against \textbf{PPO} \citep{schulman2017proximal}, which leverages a reward function learned from human preference data, and \textbf{PPO-GT}, an oracle approach that exploits the ground truth reward in controlled sentiment settings. In our sentiment-related experiments, we employ two variants of PPO-GT: a standard, unmodified implementation \cite{leandro_von_werra_2023_7790115}, as well as a custom version with reward normalization and adjusted hyperparameters for enhanced performance (these enhancements are also adopted in the conventional PPO setup with learned rewards). Lastly, the \textbf{Best of $N$} baseline is included, wherein $N$ candidate responses are sampled from the SFT model (or from Preferred-FT for dialogue tasks), and the response with the highest predicted reward—using a model trained from the preference data—is selected. While this approach can deliver strong performance and cleanly separates the quality of the reward model from PPO training, it is computationally demanding, especially as $N$ increases, since it necessitates generating $N$ completions per input at evaluation time.

\subsection{How well can DPO optimize the RLHF objective?}

The typical RLHF setup uses a reward maximization objective constrained by KL-divergence, aiming to optimize for high reward without allowing the learned policy to diverge excessively from the reference policy. Therefore, meaningful comparisons between methods should assess not only the reward obtained but also the corresponding KL divergence; gaining additional reward at the expense of significantly greater KL distance is often undesirable. Figure~\ref{fig:frontier-tldr-main} presents the reward versus KL trade-off for different approaches in the sentiment task. For each method, we conduct several independent training runs, each with distinct hyperparameters that govern the level of conservativeness (target KL $\in\{3,6,9,12\}$ for PPO, $\beta \in \{0.05,0.1,1,5\}$, $\alpha\in\{0.05,0.1,0.5,1\}$ for unlikelihood, and different random seeds for preferred-FT), resulting in a total of 22 experimental runs. At every 100 training steps up to convergence, we assess each model on a collection of test prompts, calculating both the mean reward according to the true reward function and the mean sequence-level KL\footnote{Here, sequence-level KL refers to summing the KL divergence across each timestep in the sequence.} to the reference policy $\text{KL}\left(\pi\mid \mid \pi_\text{ref}\right)$. Our experiments demonstrate that DPO defines the most favorable frontier, achieving the highest rewards at comparatively low KL divergence. This finding is striking for several reasons. Notably, while both DPO and PPO target the same objective, DPO is much more sample efficient, and its reward versus KL trade-off uniformly surpasses that of PPO. Furthermore, DPO yields a superior trade-off curve compared to PPO even in scenarios where PPO has access to true rewards (PPO-GT).

\subsection{Can DPO scale to real preference datasets?}
\label{sec:dpo-real-datasets}
We proceed to assess the effectiveness of DPO when fine-tuned on real-world preference datasets involving summarization and single-turn dialogue. In the context of summarization, it is well-documented that standard automatic metrics such as ROUGE often show weak alignment with human judgment~\citep{stiennon2022learning}. Prior studies have demonstrated that employing PPO for fine-tuning language models in accordance with human feedback can lead to improved summarization outcomes. To compare approaches, we generate completions on the test partition of the TL;DR summarization dataset and calculate each method's average win rate against reference summaries in the test set. Completions are sampled across a temperature range from 0.0 to 1.0, with corresponding win rates displayed in Figure~\ref{fig:frontier-tldr-main} (right). All methods—DPO, PPO, and Preferred-FT—start from the same GPT-J SFT checkpoint\footnote{\url{https://huggingface.co/CarperAI/openai_summarize_tldr_sft}}. Our results indicate that at a sampling temperature of 0.0, DPO attains a win rate near 61\%, outperforming PPO, which achieves approximately 57\% at its own optimal sampling temperature (also 0.0). DPO additionally surpasses the maximum win rate attained by the best of $N$ baseline. It should be emphasized that the DPO $\beta$ hyperparameter was not systematically tuned, so there may be untapped room for further improvement. Additionally, DPO demonstrates a greater degree of resilience to increases in sampling temperature compared to PPO; the latter's effectiveness can drop to the level of the original GPT-J model as temperature rises. Notably, Preferred-FT does not demonstrate a substantial advantage relative to the SFT baseline. In Section~\ref{sec:human-judgments}, we extend the comparison between DPO and PPO using direct human assessments, where DPO outputs sampled at temperature 0.25 were judged preferable 58\% of the time when compared to PPO outputs at temperature 0.

For single-turn dialogue evaluation, we examine the performance of various approaches using the portion of the Anthropic HH dataset \citep{bai2022training} test split that contains one human-assistant exchange. In the case of GPT-4, preferred completions from the test set serve as reference responses to determine the win rates of the compared approaches. Since no standard SFT model exists for this setting, our methodology begins with the pre-trained Pythia-2.8B model. We use Preferred-FT to fine-tune a reference model on selected completions, ensuring these remain within the model’s output distribution, and subsequently apply DPO for further training. Additionally, we benchmark against the top-performing completion out of 128 generated by Preferred-FT (having empirically observed the Best of $N$ metric saturates at $N=128$ for this evaluation; refer to Appendix Figure~\ref{fig:best-of-n}), as well as against a 2-shot prompt configuration of the original Pythia-2.8B model, observing that DPO consistently matches or outperforms other methods at their optimal sampling temperatures. We further assess an RLHF model trained with PPO on the Anthropic HH dataset \footnote{\url{https://huggingface.co/reciprocate/ppo_hh_pythia-6B}} provided by a recognized source \footnote{\url{https://github.com/CarperAI/trlx/tree/main/examples/hh}}; however, for this model, we are unable to identify a prompt or temperature setting that achieves results exceeding the baseline Pythia-2.8B. In light of our TL;DR experiment findings and because both PPO and the Best of 128 strategy target the same reward function, we treat the Best of 128 baseline as an approximate indicator of PPO-level performance. In summary, DPO stands out as the only method that is computationally efficient while surpassing the reference completions on the Anthropic HH dataset, and achieves results that rival or exceed those of the more computationally costly Best of 128 approach. Figure~\ref{fig:dialogue-main} additionally demonstrates that DPO attains peak performance after relatively few training steps.

\subsection{Generalization to a new input distribution}

To more rigorously assess how PPO and DPO handle distributional changes, we examine the performance of both policies—trained from our Reddit TL;DR summarization experiment—on a distinct input domain: the test split of the CNN/DailyMail dataset \citep{nallapati-etal-2016-abstractive}, comprised of news articles. We maintain the top-performing sampling temperatures previously found on the TL;DR task (0 and 0.25). Table~\ref{tab:ood} displays these outcomes. For evaluation, we calculate the win rate of GPT-4 in preference to the reference summaries from the dataset, leveraging the same GPT-4 (C) prompt as in the Reddit TL;DR evaluation, modifying only the terminology from ``forum post'' to ``news article.'' On this new input type, DPO continues to yield superior performance over PPO by a notable margin. These findings offer preliminary support that DPO can match or exceed PPO in generalizing to unseen distributions, even though DPO does not rely on the extra, unlabeled prompts from Reddit TL;DR that PPO incorporates.

\subsection{Validating GPT-4 judgments with human judgments}
\label{sec:human-judgments}

To assess the dependability of GPT-4’s evaluation capability, we conduct a human annotation study using outputs from the TL;DR summarization experiment, evaluating two distinct GPT-4 prompting strategies. The \textbf{GPT-4 (S)} (simple) prompt queries which summary better encapsulates the core information of the original post, while the \textbf{GPT-4 (C)} (concise) prompt includes an additional criterion of conciseness. The latter was selected since the simple prompt tends to bias GPT-4 towards longer and sometimes overly repetitive outputs, which are less aligned with human preferences. Complete prompt texts are provided in Appendix~\ref{app:prompts}. 

We organize three different pairwise comparisons, representing the spectrum of model performance: the highest-scoring (DPO at temperature 0.25), the lowest-performing (PPO at temperature 1.0), and an intermediate system (SFT at temperature 0.25). Each is evaluated against the PPO model sampled greedily (which performed best at its optimal temperature). Across both GPT-4 prompts, the level of agreement between GPT-4 and human annotators is on par with inter-human agreement rates, indicating GPT-4 serves as an effective stand-in for human evaluation. Because the number of human annotators was limited, multiple independent judgments were collected only for DPO and PPO-1 pairings. Notably, the \textbf{GPT-4 (C)} prompt most closely reflects human win rates, leading us to select this prompt for principal results shown in Section~\ref{sec:dpo-real-datasets}. Additional methodological information, including screenshots of the evaluation interface and a roster of participants, appears in Appendix~\ref{app:human-study}.

\section{Discussion}
Preference-based learning offers an effective and scalable approach for developing language models that are both capable and aligned with human values. In this work, we have presented DPO, a straightforward methodology for leveraging preferences to train language models, eliminating the need for reinforcement learning. Instead of adapting the preference learning problem to fit into a classical RL framework with standard RL tools, DPO provides a correspondence between policy functions and reward structures, making it possible to optimize language models in accordance with human preferences \textit{directly} by means of a cross-entropy loss—dispensing with the requirements of reinforcement learning or compromising general applicability. Remarkably, DPO attains results comparable to, or surpassing, existing RLHF strategies such as those utilizing PPO, often without significant hyperparameter adjustment. Consequently, DPO lowers the technical and computational threshold for training language models with human preference data.

\textbf{Limitations \& Future Work.} Our findings introduce several avenues for further investigation. A key question pertains to how DPO-trained policies behave when extrapolated beyond the training distribution, especially relative to those optimized using explicit reward signals. Preliminary evidence indicates that generalization capabilities of DPO are on par with models trained via PPO, but this warrants deeper empirical examination. For instance, it remains to be seen whether self-labeling mechanisms built on DPO policies can efficiently exploit unlabeled prompt data. Moreover, questions persist about the occurrence and character of reward over-optimization under direct preference optimization, and whether the slight dip in performance seen in Figure~\ref{fig:dialogue-main}-right exemplifies this phenomenon. Further, our experiments are restricted to models with up to 6B parameters, making the extension of DPO to far larger, cutting-edge model architectures a promising topic for upcoming research. In terms of evaluation methodologies, we observe that prompt choice can influence win rates reported by GPT-4, highlighting a need for systematic investigations into strategies for eliciting consistent and high-fidelity assessments from automated evaluators. Lastly, DPO may hold considerable potential for application areas beyond language model training—such as in the optimization of generative models across various modalities. 

\end{document}